经典标度律用模型参数量 $N$ 与训练数据量 $D$ 解释损失随规模的变化；面对异质语料和有限算力，还需研究如何将数据质量 $Q$ 与领域配比 $\bp$ 纳入广义标度律。本题依次涉及数据质量建模、广义标度律建模、算力约束下的模型训练优化和能力演进预测。四问环环相扣：先回答“数据如何衡量”，再回答“资源如何发挥作用、如何分配”，最后考察模型能力如何演进，逐步接近真实训练与评测场景。

针对问题一，我们将来源不同、方向与量纲各异的质量信号转为可比指标，用稳健聚合获得文档和领域质量分；把来源间的显著分歧作为人工复核线索，不将其直接判为真实质量冲突。对于领域配比，先处理份额之和固定带来的约束，再建立配比与验证损失的关系。由此得到可供后续问题使用的质量和配比接口；代理模型预测，所选配比较均匀配比的损失约低 $9.9\%$，这一数字不是大模型训练的实测收益。

针对问题二，我们以经典标度律为基线，比较质量进入损失的不同方式，并区分可由扩大规模部分补偿的质量影响与难以补偿的损失下界。在半合成数据的证据范围内，模型提示提质与扩规模可以部分替代；领域配比的作用暂作为探索性关系保留。

针对问题三，我们把训练、提质和长上下文开销放入同一预算，利用成本结构缩减求解维度，并以最优质量所处的边界状态识别资源策略的转变。现有模型显示，预算和提质成本共同改变资源的最优去向；上下文变长也会挤占可用于模型规模和数据量的算力。

针对问题四，我们区分模型类型与时间口径，先识别基准任务的饱和现象，再把损失变化映射为能力得分，分解历史前沿中的规模扩张与非规模变化，并考察算力增速放缓的情景。既有历史结果提示规模扩张是主要解释项；前沿走势保留为受预测起点与数据范围约束的条件性情景。

上述结论依赖附件覆盖的模型规模、数据来源和评测口径；尚未完成独立留出检验与当前时点重估的环节，不给出确定的收益倍数或未来分数。

\vspace{6pt}
\noindent{\heitibf 关键词：}数据质量；领域配比；广义标度律；算力约束优化；能力演进预测
